\section{Interlayer Inhibition and Adjacency Effects}
\label{sec:interlayer}

\subsection{The Mechanism}

Modern color negative films incorporate Development Inhibitor Releasing (DIR)
couplers \cite{james1977}. When development proceeds at a point in a layer, these couplers
release a soluble inhibitor that diffuses laterally within the layer and
vertically into adjacent layers, locally suppressing further development.

Two distinct and separately visible consequences follow:

\begin{enumerate}
\item \textbf{Intra-layer adjacency (acutance).} A densely developed region
suppresses development just outside its boundary, producing a light rim on the
low-density side of an edge --- the Eberhard effect. Because the suppression is
proportional to neighborhood density and subtracted locally, the net operator is
an unsharp mask applied in the density domain. This is why color negative film
has apparent sharpness exceeding what its MTF alone would predict, and why it
does not look like digital sharpening: the effect is in density, is asymmetric
in its response to edge polarity, and scales with local development activity.
\item \textbf{Inter-layer (inter-image) effect.} Inhibitor from the
green-sensitive layer suppresses development in the red and blue layers. Where a
region is strongly green, red and blue develop less, so the region prints
\emph{more} green. This is a saturation-enhancing, edge-localized, cross-channel
coupling --- and it is a substantial part of why color negative film renders
saturated colors the way it does. It cannot be reproduced by a saturation
control, because it operates on the local difference between channels rather
than on the global distance from neutral.
\end{enumerate}

\subsection{Formulation}

Let $D_k(\mathbf{p})$ be the developed density in record $k$ at position
$\mathbf{p}$ before inhibition. Inhibitor released is proportional to
development activity, which is proportional to density; the released inhibitor
diffuses with a Gaussian kernel $G_{\sigma_k}$ characteristic of the diffusion
length. The concentration of inhibitor from layer $k$ is therefore

\begin{equation}
I_k = G_{\sigma_k} * D_k,
\label{eq:inhibitor}
\end{equation}

and the suppressed density in record $c$ is

\begin{equation}
\tilde{D}_c = D_c - \sum_{k} a_{ck}\, I_k
= D_c - \sum_{k} a_{ck}\, \left(G_{\sigma_k} * D_k\right).
\label{eq:inhibited}
\end{equation}

Equation \eqref{eq:inhibited} as written contains a large DC term: it lowers
mean density everywhere. That component is real, but it is \emph{already
present in the measured characteristic curve} to which $\theta$ was fitted ---
sensitometric measurement of a real film necessarily includes its own
inhibition. Applying \eqref{eq:inhibited} directly would therefore double-count.

The correct treatment is to decompose the inhibitor into its local mean and its
spatial residual. Writing the highpass operator as

\begin{equation}
H_{\sigma}[D] \;=\; D - G_{\sigma} * D,
\end{equation}

we have $G_\sigma * D_k = D_k - H_{\sigma_k}[D_k]$, so \eqref{eq:inhibited}
becomes

\begin{equation}
\tilde{D}_c = \underbrace{D_c - \sum_k a_{ck} D_k}_{\text{pointwise; absorbed into } \theta}
\;+\; \sum_k a_{ck}\, H_{\sigma_k}[D_k].
\end{equation}

The runtime operator is therefore the residual term alone:

\begin{equation}
\boxed{\;\tilde{D}_c \;=\; D_c \;+\; \Lambda_c \sum_{k \in \{R,G,B\}} a_{ck}\, H_{\sigma_k}\!\left[D_k\right]\;}
\label{eq:interlayer}
\end{equation}

where $\Lambda_c$ is an activity weight, discussed next. This form has the
essential property that it is \emph{mean-preserving}: $H_\sigma$ annihilates
constants, so a uniform patch is unchanged and no double-counting occurs. All of
the visible effect --- edge rims and local saturation --- is retained.

\subsection{Activity Weighting}

Inhibitor release is proportional to development activity, not to density. In
the toe, almost nothing is developing and almost no inhibitor is released; in
the shoulder, development has run to completion and further release does
nothing. Both are captured by the point gamma of
\eqref{eq:pointgamma}, normalized:

\begin{equation}
\Lambda_c(\mathbf{p}) \;=\; \left(\frac{\gamma_c(x_c(\mathbf{p}))}{\gamma_c}\right)^{\!\mu},
\qquad \mu \approx 1.0.
\label{eq:activityweight}
\end{equation}

Consequently the adjacency effect is strongest in the mid-scale and vanishes in
both extremes --- exactly the observed behavior, and a meaningful difference
from an unsharp mask, which sharpens deep shadows and blown highlights with
equal enthusiasm.

Because $\gamma_c(x)$ is already computed by \eqref{eq:pointgamma} for the grain
model, $\Lambda_c$ costs nothing additional.

\subsection{Multi-Scale Diffusion}

A single Gaussian is a poor model of diffusion in a layered gelatin structure,
which exhibits both a short-range component within the layer and a long-range
component through the interlayer. \EM{} uses a two-scale kernel:

\begin{equation}
H^{(2)}[D] = w_1\,H_{\sigma_1}[D] + w_2\,H_{\sigma_2}[D],
\qquad \sigma_2 \approx 5\sigma_1,
\label{eq:twoscale}
\end{equation}

with $w_1 + w_2 = 1$. The short scale $\sigma_1 \approx 1.2\,\mu\mathrm{m}$
produces acutance; the long scale $\sigma_2 \approx 6\,\mu\mathrm{m}$ produces
the broad-area inter-image effect. Both are specified in \emph{micrometres at
the film plane} and converted to pixels through the format scaling of
Section~\ref{sec:formatscale}, so that the effect has a fixed physical size
regardless of output resolution.

\subsection{The Coupling Matrix}

The matrix $\mathbf{A} = [a_{ck}]$ is the stock's inhibition signature. Its
structure is constrained by physics:

\begin{itemize}
\item Diagonal entries $a_{cc} > 0$: self-inhibition, producing acutance.
\item Off-diagonal entries $a_{ck} > 0$ for $k \neq c$: cross-layer inhibition.
\item Vertical adjacency matters. In the standard layer order (blue on top,
then green, then red), green couples strongly to both neighbors, while
red-to-blue coupling is weak because the layers are not adjacent and the yellow
filter layer intervenes. The matrix is therefore approximately tridiagonal in
the layer ordering, not symmetric in the RGB ordering.
\end{itemize}

A representative shipping matrix for a modern color negative stock, in layer
order $(B, G, R)$ and re-expressed in RGB indexing, is

\begin{equation}
\mathbf{A} =
\begin{pmatrix}
0.42 & 0.26 & 0.05 \\
0.19 & 0.48 & 0.19 \\
0.05 & 0.29 & 0.38
\end{pmatrix},
\label{eq:dirmatrix}
\end{equation}

where row index is the affected record $c$ and column index the releasing
record $k$. Note the small $R\!\leftrightarrow\!B$ terms and the strong
coupling of both to $G$, reflecting the physical stack.

A monochrome stock has a scalar $a$ and no cross terms. A transparency stock has
substantially weaker inhibition overall (reversal films use fewer DIR couplers),
with all entries scaled by roughly $0.4$.

\subsection{Implementation}

The operator is separable and cheap. The pass computes two Gaussian blurs of the
density image at $\sigma_1$ and $\sigma_2$, forms the two highpass residuals,
applies \eqref{eq:twoscale}, multiplies by $\mathbf{A}$ and by $\Lambda$, and
adds. Blurs use the separable recursive filter of
Section~\ref{sec:recursive-blur}, so cost is independent of $\sigma$.

\begin{lstlisting}[language=Metal, caption={Interlayer inhibition combine pass.}]
kernel void interlayer_combine(
    texture2d<float, access::read>  density   [[texture(0)]],
    texture2d<float, access::read>  blur1     [[texture(1)]],
    texture2d<float, access::read>  blur2     [[texture(2)]],
    texture2d<float, access::read>  gammaMap  [[texture(3)]],
    texture2d<float, access::write> outTex    [[texture(4)]],
    constant InterlayerParams&      P         [[buffer(0)]],
    uint2 gid [[thread_position_in_grid]])
{
    if (gid.x >= outTex.get_width() || gid.y >= outTex.get_height()) return;

    float3 D  = density.read(gid).rgb;
    float3 b1 = blur1.read(gid).rgb;
    float3 b2 = blur2.read(gid).rgb;

    // Two-scale highpass residual, eq. (two-scale)
    float3 H = P.w1 * (D - b1) + P.w2 * (D - b2);

    // Coupling matrix A applied across records
    float3 coupled = P.A * H;             // float3x3 * float3

    // Activity weight from normalized point gamma
    float3 lambda = pow(max(gammaMap.read(gid).rgb, 1e-4f), P.mu);

    outTex.write(float4(D + P.strength * lambda * coupled, 1.0f), gid);
}
\end{lstlisting}

\subsection{User Exposure of the Effect}

Parameter honesty forbids labelling this control ``sharpness.'' The exposed
controls are:

\begin{itemize}
\item \textbf{Coupler activity} ($0$--$200\%$, default $100\%$), scaling
$\mathbf{A}$ uniformly. This is the strength of the DIR effect.
\item \textbf{Agitation} (from Section~\ref{sec:development}), which
\emph{inversely} modulates $\sigma_2$ and $w_2$: less agitation, longer-range
inhibition, stronger broad-area effects. This is the single control that most
rewards a user who knows darkroom practice.
\end{itemize}

Intra-layer and inter-layer strength are not separately exposed in v1.0, because
they are not separately controllable in a real process. They are, however,
separately editable in the stock profile format, which is how a future
user-authored-profile feature would surface them.
